\section{Data Sources}

This project uses three main kinds of data: tree cadastre records [insert citation], Sentinel-2 satellite imagery [insert citation], and OpenStreetMap (OSM) urban features [insert citation]. Together, these sources support the project goal of mapping urban tree canopy from freely available geospatial data. They have already been gathered, cleaned, and stored in reusable formats using Jupyter notebooks [insert Nina GitHub citation], so they can later be used for exploratory data analysis and modelling.
\\
The data sources and their cleaning processes are described in the following subchapters.

\subsection{Tree cadastre data}

Four cities are included: Vienna [insert citation], Paris [insert citation], Barcelona [insert citation], and Hamburg [insert citation]. These datasets were loaded from local raw-data files in different formats. The cadastre sources are important because they provide point-based information about known trees. They are the main direct source of tree locations in this project.
\\
The data preparation follows several steps. First, duplicates are checked, next, columns that are not needed for modelling are removed, while the most useful variables are retained, such as geometry, tree species, height, circumference, crown diameter, and a city identifier. The notebook also standardizes coordinate systems and column names to make the datasets comparable across cities.Missing values are handled differently depending on the city. Vienna, Paris, and Barcelona do not contain missing values in the retained fields after cleaning. In the Hamburg Port Authority data, crown diameter and circumference still contain missing values after the first cleaning step. Instead of dropping these rows, the missing values are imputed using the median within the same species, with the overall median used as a fallback. Unknown species names are replaced with the placeholder value unknown.
\\
After cleaning, the cadastre data are saved in GeoPackage and also Parquet, for compact storage. Barcelona is saved only as Parquet, since its processed version is not stored as a GeoPandas geometry dataset in the same way as the other city datasets.

\subsection{Sentinel-2 satellite imagery}

For Sentinel-2 data collection, we first considered the Copernicus API [insert citation], but then switched to Google Earth Engine[insert citation] and created monthly median composites from Sentinel-2 images filtered by cloud coverage, which makes the process more reproducible and reduces cloud-related problems compared with manual image selection in the previous project. Google Earth Engine also allowed us to restrict the data retrieval to the exact area used in the tree cadastre data.
\\
The project uses seasonal satellite imagery for spring, summer, and autumn, mainly April, August, and November. These months were chosen to capture different stages of vegetation development. In Paris, October was used instead of November because no suitable November image was available.The image collection is filtered by area of interest, date, and cloud coverage (CLOUDY-PIXEL-PERCENTAGE < 20). A stricter 10 percent threshold was tested first, but it did not provide usable images for all months. From Sentinel-2, the workflow keeps four bands: B02, B03, B04, and B08, [insert citation why those] corresponding to blue, green, red, and near-infrared. In addition, three vegetation indices are computed: NDVI, EVI, and SAVI. [insert citation why those] 
\\
The result is a cleaned multi-season spectral stack for each city.

\subsection{OpenStreetMap urban features}
We also collect urban context layers from OpenStreetMap (OSM) [insert citation]. These data come from Geofabrik downloads[insert citation] in GeoPackage format. Separate regional files are used for Austria, Île-de-France, Catalonia, and Hamburg. From each source, the notebook reads four thematic layers: roads, buildings, land use, and water.
\\
After loading the OSM layers and converting them to EPSG:4326, they are clipped to the same city bounding boxes derived from the cleaned tree-cadastre data. This keeps the cadastre data, Sentinel-2 imagery, and OSM layers spatially aligned within the same study areas. At this stage, the OSM data are only lightly cleaned. We also check geometry validity and empty geometries, but no strong filtering of feature classes is applied yet to preserve flexibility for later modelling. 
\\
The cleaned layers are saved for reuse in GeoPackage and Parquet format and are prepared the exploratory data analysis.



